\documentclass[11pt,a4paper]{article}

\usepackage[T1]{fontenc}
\usepackage[utf8]{inputenc}
\usepackage{lmodern}
\usepackage{microtype}
\usepackage[a4paper,margin=24mm,headheight=22pt]{geometry}
\usepackage[dvipsnames,table]{xcolor}
\usepackage{enumitem}
\usepackage{booktabs}
\usepackage{tabularx}
\usepackage{listings}
\usepackage{fancyhdr}
\usepackage{titlesec}
\usepackage{hyperref}
\usepackage{bookmark}

\definecolor{GuideNavy}{HTML}{17324D}
\definecolor{GuideTeal}{HTML}{167A7B}
\definecolor{GuidePale}{HTML}{EEF6F6}
\definecolor{GuideWarn}{HTML}{8B4A16}
\definecolor{GuideWarnPale}{HTML}{FFF4E8}
\definecolor{GuideCode}{HTML}{F3F5F7}
\definecolor{GuideRule}{HTML}{C9D4DD}

\hypersetup{
  colorlinks=true,
  linkcolor=GuideTeal,
  urlcolor=GuideTeal,
  pdftitle={Austrian ICD-10 Educational Prototype -- User Guide},
  pdfauthor={bsc-thesis-icd10 project},
  pdfsubject={Installation, use, maintenance, and troubleshooting of the prototype},
  pdfkeywords={ICD-10, Austria, education, prototype, user guide}
}

\pagestyle{fancy}
\fancyhf{}
\fancyhead[L]{\small\color{GuideNavy}Austrian ICD-10 Educational Prototype}
\fancyhead[R]{\small\color{GuideNavy}User Guide}
\fancyfoot[C]{\small\color{GuideNavy}\thepage}
\renewcommand{\headrulewidth}{0.4pt}
\renewcommand{\headrule}{\hbox to\headwidth{\color{GuideRule}\leaders\hrule height \headrulewidth\hfill}}

\titleformat{\section}
  {\Large\bfseries\color{GuideNavy}}{\thesection}{0.7em}{}
\titleformat{\subsection}
  {\large\bfseries\color{GuideTeal}}{\thesubsection}{0.7em}{}
\titlespacing*{\section}{0pt}{2.2ex plus 0.5ex}{1.0ex}
\titlespacing*{\subsection}{0pt}{1.8ex plus 0.4ex}{0.7ex}

\setlist[itemize]{leftmargin=1.5em,itemsep=0.35em,topsep=0.5em}
\setlist[enumerate]{leftmargin=1.7em,itemsep=0.45em,topsep=0.5em}
\setlength{\parindent}{0pt}
\setlength{\parskip}{0.55em}
\setlength{\emergencystretch}{2em}

\lstdefinestyle{terminal}{
  basicstyle=\ttfamily\small,
  backgroundcolor=\color{GuideCode},
  frame=single,
  rulecolor=\color{GuideRule},
  framesep=7pt,
  xleftmargin=2pt,
  xrightmargin=2pt,
  breaklines=true,
  breakatwhitespace=false,
  columns=fullflexible,
  keepspaces=true,
  showstringspaces=false,
  aboveskip=0.8em,
  belowskip=0.8em
}

\newcommand{\productname}{Austrian ICD-10 Educational Prototype}
\newcommand{\infobox}[2]{%
  \par\medskip
  \noindent\fcolorbox{GuideTeal}{GuidePale}{%
    \begin{minipage}{0.94\linewidth}
      \textbf{\color{GuideNavy}#1}\par\smallskip
      #2
    \end{minipage}%
  }
  \par\medskip
}
\newcommand{\warningbox}[2]{%
  \par\medskip
  \noindent\fcolorbox{GuideWarn}{GuideWarnPale}{%
    \begin{minipage}{0.94\linewidth}
      \textbf{\color{GuideWarn}#1}\par\smallskip
      #2
    \end{minipage}%
  }
  \par\medskip
}

\begin{document}
\hypersetup{pageanchor=false}

\begin{titlepage}
  \thispagestyle{empty}
  \vspace*{23mm}
  {\color{GuideTeal}\rule{\textwidth}{2.2pt}}\par
  \vspace{13mm}
  {\Huge\bfseries\color{GuideNavy} Austrian ICD-10\\[0.18em]
    Educational Prototype\par}
  \vspace{5mm}
  {\LARGE\color{GuideTeal} User Guide\par}
  \vspace{12mm}
  {\large Installation, learner workflow, maintenance, and troubleshooting\par}
  \vfill
  \infobox{At a glance}{
    This guide is for people who want to install and use the working thesis
    prototype. The quickest route uses Docker Compose and the project's
    published container images. No local PHP, MySQL, Node.js, Python, or
    Composer installation is required.
  }
  \vfill
  \begin{tabular}{@{}ll@{}}
    \textbf{Guide revision} & 0.3 \\
    \textbf{Last updated} & 9 August 2026 \\
    \textbf{Prototype data} & BMASGPK 2026 working baseline (6 patients, 25 questions) \\
    \textbf{Repository} & \href{https://github.com/junomarx/bsc-thesis-icd10}{github.com/junomarx/bsc-thesis-icd10}
  \end{tabular}\par
  \vspace{12mm}
  {\color{GuideTeal}\rule{\textwidth}{2.2pt}}\par
\end{titlepage}

\hypersetup{pageanchor=true}
\pagenumbering{roman}
\tableofcontents
\clearpage
\pagenumbering{arabic}

\section{Quick introduction}

\productname{} is a bachelor-thesis software artefact for practising a
small, explicitly defined subset of Austrian ICD-10 coding. It presents six
synthetic teaching patients, each containing several independent coding
questions; the learner answers one question at a time by selecting a
supported code (or \textbf{None of the above}), and each answer is
evaluated with deterministic rules. Feedback is returned as
\textbf{Correct}, \textbf{Suboptimal}, or \textbf{Incorrect}, together with
an explanation and, where applicable, a suggested improvement. The
interface is available in English and German.

The prototype is deliberately small. It is intended to demonstrate and test
the project's rule-based educational approach, not to provide a complete
coding platform.

\warningbox{Educational use only}{
  The patients and questions are synthetic. The prototype does not diagnose
  patients, provide clinical decision support, perform official coding, or
  support reporting or reimbursement. Do not enter real patient data.
}

\section{Installation}

\subsection{What you need}

Before starting, install and start the following:

\begin{itemize}
  \item \textbf{Docker Desktop} on macOS or Windows, or Docker Engine on
        Linux;
  \item \textbf{Docker Compose v2}, available through the
        \texttt{docker compose} command;
  \item \textbf{Git} to download and update the repository;
  \item a modern web browser; and
  \item an internet connection for the initial download of the repository and
        container images.
\end{itemize}

Check Docker before continuing:

\begin{lstlisting}[style=terminal]
docker --version
docker compose version
\end{lstlisting}

Both commands should print version information. If Docker reports that it
cannot connect to the daemon, open Docker Desktop or start the Docker service.

\subsection{Install and start with the published images}

\warningbox{The published images currently predate this guide's learner workflow}{
  \texttt{docker compose pull} fetches the last images published by the
  project's CI pipeline, which has not been re-run since the six-patient/
  25-question model described in Section~\ref{sec:learner-workflow} was
  implemented. Following the two commands below as written will currently
  start the \emph{previous} case-based prototype, not the one this guide
  describes. Until the images are republished, use
  Section~\ref{sec:local-build}'s native local build instead ---
  \texttt{docker compose build bootstrap app} followed by
  \texttt{docker compose up -d --wait app} --- which builds directly from
  the checked-out source and is guaranteed to match this guide. This
  warning will be removed once a fresh publish confirms the pulled images
  are current again.
}

Open a terminal and run:

\begin{lstlisting}[style=terminal]
git clone https://github.com/junomarx/bsc-thesis-icd10.git
cd bsc-thesis-icd10
docker compose pull
docker compose up -d --wait app
\end{lstlisting}

The first start downloads several images and may take a few minutes. Compose
then starts MySQL, loads the versioned prototype baseline, and starts the web
application in the correct order.

The project-owned images are published for both \texttt{linux/amd64}
(Intel/AMD computers) and \texttt{linux/arm64} (including Apple Silicon).
Docker selects the matching variant automatically.

\infobox{If an older image tag has no Apple Silicon variant}{
  If \texttt{docker compose pull} reports
  \texttt{no matching manifest for linux/arm64/v8}, the registry still has an
  older, AMD64-only publication. Use the native source-build procedure in
  Section~\ref{sec:local-build}; it works without emulation and does not depend
  on the registry containing the current multi-platform release.
}

Open \url{http://localhost:5860} in a browser.

Optionally verify the HTTP endpoint from a terminal:

\begin{lstlisting}[style=terminal]
curl http://127.0.0.1:5860/api/health
\end{lstlisting}

The expected response is:

\begin{lstlisting}[style=terminal]
{"status":"ok"}
\end{lstlisting}

\infobox{What was installed?}{
  Docker created containers for the database, one-time baseline loader, and
  application, plus a named database volume. The source checkout remains a
  normal folder and can be updated with Git. Learner attempts are not stored.
}

\subsection{Use a different local port}
\label{sec:port-fallback}

Port 5860 is the default. If it is already in use, stop the stack and choose
another port. On macOS or Linux:

\begin{lstlisting}[style=terminal]
docker compose down
APP_HTTP_PORT=8081 docker compose up -d --wait app
\end{lstlisting}

In Windows PowerShell:

\begin{lstlisting}[style=terminal]
docker compose down
$env:APP_HTTP_PORT = "8081"
docker compose up -d --wait app
\end{lstlisting}

Then open \url{http://localhost:8081}. The selected environment value remains
in the current terminal session; clear or change it before a later start if
you want to return to port 5860.

\subsection{Build locally instead of pulling images}
\label{sec:local-build}

If a published image is unavailable for the host architecture, or if you want
to run the exact source in the checkout, build the two images needed for
normal use locally:

\begin{lstlisting}[style=terminal]
docker compose build bootstrap app
docker compose up -d --wait app
\end{lstlisting}

This takes longer on the first run because Docker builds the frontend,
backend, and baseline-loader images. It still does not require the language
toolchains to be installed directly on the host. On Apple Silicon, the
resulting \texttt{app} and \texttt{bootstrap} images are native ARM64 images.

\section{Using the prototype}
\label{sec:learner-workflow}

\subsection{Switch language}

An \textbf{EN | DE} switch sits in the top-right corner of every screen.
The choice is remembered in the browser for later visits, and defaults to
German or English automatically based on the browser's own language
setting on first visit. Switching language translates all interface text,
the patient summaries and question prompts, and the feedback explanation.
The displayed ICD-10 code designations come from the single-language
Austrian source catalogue and are shown in the language of that source
regardless of the interface language: German designations are shown as
authored; the prototype supplies English titles for these instead when the
interface is set to English, since no official English designation exists
in the source data.

\subsection{How this works: orientation}

Below the page header, an expandable \textbf{How this works} panel states
the prototype's educational purpose and non-clinical boundary, summarizes
the choose-patient $\rightarrow$ answer-question $\rightarrow$
receive-feedback workflow, and explains the three feedback classes with the
same icon and colour used later in the actual feedback. It is open by
default and can be collapsed; collapsing it does not lose any progress.

\subsection{Choose a patient}

Below the orientation panel, the six synthetic patients appear as a grid of
equally sized cards. Each card shows the patient's display name, age, sex,
a neutral complexity label (\textit{Foundational} or \textit{Involved}),
and its question count (3, 5, or 6 depending on the patient), followed by a
one-sentence summary of that patient's record. Difficulty is informational
only --- no patient is locked, and any patient may be selected in any
order.

A patient already fully answered in this browser session shows a
\textit{Completed} badge on its card, and a one-line summary above the
grid tracks overall progress ("\textit{n} of 6 patients completed''). This
record lives only for the current browser session (cleared when the
browser tab/session ends, not just on page reload) --- nothing is sent to
or stored by the server. A \textbf{Reset progress} link appears next to the
summary whenever at least one patient has been completed, and clears every
completion mark after a confirmation prompt. Select a patient to begin.

\subsection{Review the patient record}

Choosing a patient opens its first question directly. A \textbf{Show
patient info} button reveals a collapsible panel with the patient's
identity, a general summary, and its documented context items (conditions,
history, current findings), each labelled by information type --- for
example, whether a detail comes from a documented record versus a current
examination finding, which matters for questions about a patient who
cannot provide their own history. This panel can be opened and closed at
any point without losing the current question's state, and remains
reachable throughout the patient's questions via the same button.

\subsection{Answer a question}

Each question shows its position (``Question 2 of 5''), a matching visual
progress bar, the patient panel, the question prompt, and a list of
answer options: several supported ICD-10 codes with their designation,
plus a dedicated \textbf{None of the above} option. Select one option, then
choose \textbf{Submit answer}; the button remains disabled until an option
is selected. Question order is shuffled independently each time a patient
is started or replayed, so the same patient may present its questions in a
different sequence on a later attempt --- this never changes which
questions or options exist, only their order.

Only the options shown are offered; free-text or arbitrary code submission
is not part of the learner interface. An \textbf{Exit to patient list}
link is available at any point during a question, in case you want to
return to the patient list before finishing --- a confirmation is required,
since the current, unfinished attempt on that patient is not saved.

\subsection{Interpret the feedback}

Each result is shown with a matching icon and colour alongside its text
label, so the outcome does not depend on colour perception alone.

\begin{tabularx}{\textwidth}{@{}>{\bfseries}p{28mm}X@{}}
\toprule
Result & Meaning in this prototype \\
\midrule
Correct & The submitted code is a declared acceptable response for the
question. \\
Suboptimal & The answer is accepted within the implemented rules, but the
documented facts support a more specific code. A suggested improvement is
shown. \\
Incorrect & A hard rule was violated, for example a status restriction,
required coding depth, conflict with a documented fact, or a wrong
temporal/status context --- or, for \textbf{None of the above}, a declared
acceptable code was in fact among the displayed options. \\
Not evaluated & The request could not be classified because a required
question or code relationship was unavailable. This is normally a data or
application problem, not an expected learner outcome. \\
\bottomrule
\end{tabularx}

The explanation beneath the result states the rule-relevant reason, in the
selected interface language. A collapsed \textbf{Technical details}
section is available beneath it, showing the determining rule and
criterion for anyone who wants to inspect exactly which rule fired ---
useful for demonstration, not required reading for normal use. Once
submitted, an answer is locked for the rest of that attempt; choose
\textbf{Next question}, or \textbf{Review patient} after the last question
of that patient.

\subsection{Review a completed patient}

After the last question, a review screen shows the patient's name with a
completion acknowledgement, raw counts of correct/suboptimal/incorrect
results (never a weighted score), and a list of every question with its
result. From here, choose \textbf{Play again} to reattempt the same
patient with its questions reshuffled, or \textbf{Choose another patient}
to return to the patient list.

\section{Starting, stopping, and updating}

Run all commands in the repository folder containing
\texttt{docker-compose.yml}.

\subsection{Check status and logs}

\begin{lstlisting}[style=terminal]
docker compose ps
docker compose logs app
docker compose logs db
docker compose logs bootstrap
\end{lstlisting}

Add \texttt{-f} to a log command to follow new output, for example
\texttt{docker compose logs -f app}. Press Ctrl+C to stop following the log;
this does not stop the container.

\subsection{Stop and restart}

Stop the containers while keeping the database volume:

\begin{lstlisting}[style=terminal]
docker compose down
\end{lstlisting}

Start the prototype again later:

\begin{lstlisting}[style=terminal]
docker compose up -d --wait app
\end{lstlisting}

\subsection{Update to the newest repository and images}

\begin{lstlisting}[style=terminal]
git pull --ff-only
docker compose pull
docker compose up -d --wait app
\end{lstlisting}

The bootstrap service checks the versioned baseline during startup and is
safe to run again when the baseline is unchanged.

\subsection{Reset the prototype data}
\label{sec:reset-data}

\warningbox{This removes the database volume}{
  The following command deletes the local Compose database volume. In the
  current prototype that volume contains only the reproducible, versioned
  baseline---not learner accounts or attempt histories. Do not reuse this
  command unchanged if persistent user data is added in a future version.
}

\begin{lstlisting}[style=terminal]
docker compose down -v
docker compose up -d --wait app
\end{lstlisting}

Resetting is useful after an incompatible \texttt{mysql:latest} upgrade or
when you need a clean copy of the baseline.

\section{Optional: run the automated tests}

\warningbox{Currently broken, pending an in-progress test-suite update}{
  The committed PHPUnit suite (\texttt{app/tests/}) still targets the
  prototype's previous, single-question-per-case data model and has not
  yet been rewritten for the current six-patient/25-question model
  described in this guide. Running the commands below currently reports
  most unit tests erroring, not passing. This is a known, tracked gap
  (see \path{docs/CHANGELOG.md}), not a sign that something in your local
  setup is broken. The application itself is unaffected --- it does not
  depend on this test suite to run.
}

The application does not start test tooling during normal use. To run the
containerized unit, integration, and browser-driven test suite:

\begin{lstlisting}[style=terminal]
docker compose --profile test up -d --wait selenium
docker compose --profile test run --rm test
\end{lstlisting}

If the test image also needs to be built from the checkout, run this once
before the commands above:

\begin{lstlisting}[style=terminal]
docker compose --profile test build test
\end{lstlisting}

Once the test-suite update above lands, a successful run will end with
PHPUnit reporting no failures or errors; until then, expect errors as
described in the warning above. The browser tests use a separate Selenium
container and may take longer on the first run.

Stop the application and test services when finished:

\begin{lstlisting}[style=terminal]
docker compose --profile test down
\end{lstlisting}

The database volume is retained unless \texttt{-v} is added.

\section{Troubleshooting}

\subsection{Docker cannot connect}

If a command reports that the Docker daemon is unavailable, start Docker
Desktop or the Linux Docker service and retry \texttt{docker compose ps}.

\subsection{Port 5860 is already allocated}

Use the different-port procedure in Section~\ref{sec:port-fallback}. If a
previous copy of this project is still running, \texttt{docker compose down}
from its repository folder may release the port.

\subsection{The database is unhealthy after an update}

The project deliberately uses \texttt{mysql:latest}. MySQL may reject an
existing data volume after a major-version jump. Inspect the database log:

\begin{lstlisting}[style=terminal]
docker compose logs db
\end{lstlisting}

If the log reports an unsupported in-place upgrade, use the reset procedure
in Section~\ref{sec:reset-data}. This deletes and reconstructs the local
baseline database.

\subsection{The page does not load or patients do not appear}

Check the service state and health endpoint:

\begin{lstlisting}[style=terminal]
docker compose ps
curl http://127.0.0.1:5860/api/health
docker compose logs app
docker compose logs bootstrap
\end{lstlisting}

If you selected a different port, replace 5860. A failed bootstrap prevents
the application from starting because the baseline is a required dependency;
the bootstrap log contains the cause.

\subsection{Pull reports ``no matching manifest''}

This error means the selected registry tag does not contain an image for the
computer's architecture. It was observed on Apple Silicon with the project's
original AMD64-only publication. The publishing workflow now produces both
AMD64 and ARM64 variants. Update the checkout and retry the pull if a custom,
cached, or historical tag still exposes the earlier single-architecture index.

Run the native fallback from Section~\ref{sec:local-build}:

\begin{lstlisting}[style=terminal]
docker compose build bootstrap app
docker compose up -d --wait app
\end{lstlisting}

Do not add a permanent \texttt{platform: linux/amd64} override merely to hide
the error; that would force slower emulation on Apple Silicon after native
images are available.

\subsection{The browser tests cannot start}

Confirm that Docker has enough memory available and that ports 4444 and 7900
are not already occupied. The normal learner-facing application does not need
these ports; they belong only to the optional Selenium service.

\section{Data, privacy, and scope}

\begin{itemize}
  \item All bundled patients and questions are synthetic teaching content.
  \item The current interface has no account system and does not persist
        learner submissions or attempt history anywhere server-side. The
        per-patient \textit{Completed} marks described in
        Section~\ref{sec:learner-workflow} live only in the browser, only
        for the current session, and are never sent to the server.
  \item The database volume holds the versioned prototype catalogue,
        patient, question, rule-support, and reference-response baseline.
  \item The implemented catalogue is a small, traceable subset of the Austrian
        ICD-10 BMASGPK 2026 material, not the entire catalogue.
  \item Results express the prototype's declared rules and question
        relationships; they are not medical or official coding advice.
\end{itemize}

\section{Further project documentation}

This guide covers installation and day-to-day use. The repository contains
the following deeper references:

\begin{itemize}
  \item \path{docs/IMPLEMENTATION_SPECIFICATION.md} --- exact implemented
        schema, API, build, and deployment contracts;
  \item \path{docs/DEVELOPMENT_DOCUMENTATION.md} --- technology and design
        rationale;
  \item \path{docs/REQUIREMENTS_TRACEABILITY.md} --- requirement-to-evidence
        mapping;
  \item \path{docs/CHANGELOG.md} --- dated implementation and documentation
        history; and
  \item \path{HANDOFF.md} --- current project snapshot and remaining work.
\end{itemize}

The editable source for this document is \path{docs/USER_GUIDE.tex}; the
compiled edition is \path{docs/USER_GUIDE.pdf}. Both are maintained in the
repository so the readable PDF and its source can be reviewed together.

\appendix
\section{Command summary}

\begin{tabularx}{\textwidth}{@{}>{\raggedright\arraybackslash\ttfamily}p{69mm}X@{}}
\toprule
\normalfont\bfseries Command & \textbf{Purpose} \\
\midrule
docker compose pull & Download published images. \\
docker compose up -d --wait app & Start the application and its dependencies. \\
docker compose ps & Show service state. \\
docker compose logs app & Show application logs. \\
docker compose down & Stop containers; keep the database volume. \\
docker compose down -v & Stop containers and delete the database volume. \\
docker compose build bootstrap app & Build the normal-use images natively from the local checkout. \\
docker compose --profile test run --rm test & Run the full automated test suite after Selenium is started. \\
\bottomrule
\end{tabularx}

\end{document}
